\documentclass[12pt]{article}

\usepackage[english]{babel}
\usepackage[utf8x]{inputenc}
\usepackage{pdfpages}
\usepackage{lastpage} % Required to determine the last page for the footer
\usepackage{extramarks} % Required for headers and footers
\usepackage{graphicx} % Required to insert images
\usepackage{listings} % Required for insertion of code
\usepackage{courier} % Required for the courier font
\usepackage{color}
\usepackage{grffile}
\usepackage{float}
\usepackage{seqsplit}

\usepackage[a4paper, total={6in, 8in}]{geometry}

% Margins
\topmargin=-0.45in
\evensidemargin=0in
\oddsidemargin=0in
\textwidth=6.5in
\textheight=9.0in
\headsep=0.25in
\fboxsep=0mm%padding thickness
\fboxrule=2pt%border thickness

\linespread{1.1} % Line spacing

\newcommand{\Title}{Deployment guidelines document} % Assignment title
\newcommand{\Class}{Cos\ 301} % Course/class
\newcommand{\pd}{Post-Doctoral}
\newcommand{\ssr}{Soft\color{green}{Serve }\color{black}}
\newcommand{\version}{1.0 (Final)}
\newcommand{\iteration}{6}
\newcommand{\client}{Ms. Cathy Sandis (UP DRIS)}
\newcommand{\supervisor}{Prof. Stefan Gruner (SSFM Group)}
\newcommand{\project}{Post-Doctoral management System}
\newcommand{\repo}{https://github.com/mox1990/Project-Postdoc.git}

\begin{document}


\vspace{4em}

\begin{center}%

\begin{figure}[ht!]
\centering
\includegraphics{../Images_Docs/logo.png}
\end{figure}
\LARGE \bf \Class \\[0.25em]
\LARGE \bf \project \\[1em]
\LARGE \bf \Title \\[0.25em]
\large \bf \today\\
\bf Version \version\\
\bf Iteration \iteration\\[0.5em]
\Large \bf Prepared for \\Client: \client\\Supervisor: \supervisor
\Large \\\bf by \\
\Large {\bf \ssr Group }\\[0.5em]
\LARGE {\bf Group members}\\[0.25em]
\large
Kgothatso Phatedi Alfred Ngako (12236731) \\[0.5em]
Tokologo “Carlo” Machaba (12078027) \\[0.5em]
Mathys Ellis (12019837) \\[8em]

\end{center}%

%\newpage
%{\LARGE \bf Change log}\\[2em]

\begin{center}
\begin{tabular}{|l|p{1.4cm}|p{8cm}|p{2.8cm}|}
\hline
\multicolumn{4}{|c|}{\bf Change log} \\
\hline
 Date & Version & Description &  Person \\
\hline
10/02/2014 & v 0.0 & Created deployment guide & Alfred Ngako \\
\hline

%\end{tabbing}
\end{tabular}
\end{center}
\newpage
\tableofcontents

\listoffigures
\newpage
\section{Document description:}
This document is written to descride the deployment process required to install the Postdoctoral Management system on a server.
\subsection{Document purpose:}
\vspace{0.2in}
This document serves as a general TODO list since there existe multiple platforms and environments in which a server can be run. It has been written under the assumption that the reader has basic knowledge of all the components being used in the system. A variation of steps will be provided for the Ubuntu and Windows operating systems. The differences in installing the software is handled in the links provided. In addition, the document will provide two different implementations for developer deployment and production deployment.
\vspace{0.2in}
\section{System Requirements}
The system will require the following to be achieve successful deployment.
\begin{itemize}
\item 1 GHz 32-bit (x86) or 64-bit (x64) processor
\item At least 1GB of system memory. Recommended 2GB for Windows platforms, 1GB otherwise.
\item 10GB avaible disk space
\item Java 8 sdk
\item Glassfish server 4.0
\item MySQL 5.6
\item Netbeans 8 (for developer deployment)
\end{itemize}
The deployer also needs to have the projects latest repo downloaded on their system. the repo can be found here \textbf{\seqsplit{https://github.com/mox1990/Project-Postdoc.git}}.
\section{Prequiste software}
Before the system can be deployed a number of software requires to be installed and correctly configured. The following a number of steps to be followed to achieve this goal.
\subsection{Java 8 SDK}
This is the primary programming langauge which the system is written in and it provides both the compile and runtime environment of the system.
\subsubsection{Installation}
The software can be downloaded from oracle on this link \textbf{\seqsplit{http://www.oracle.com/technetwork/java/javase/downloads/jdk8-downloads-2133151.html}} \\
\\
An overview of the installation can be found here \textbf{\seqsplit{http://docs.oracle.com/javase/8/docs/technotes/guides/install/install\_ overview.html}}
\subsubsection{Configuration}
No particular configuration is required. As long as this sdk is the default version on the machine/IDE being used for deployment, all aspects should run as expected.
\subsection{Glassfish Server 4.0}
This is an application server which runs the Java EE 7 application.
\subsubsection{Installation}
All the required steps for installation are found here \textbf{\seqsplit{https://glassfish.java.net/getstarted.html}}.
\subsubsection{Configuration}
The application requires the use of a MySQL database as a resource. To do this, the JDBC resources (which can be found under \textbf{\seqsplit{~/PostDoctoralManagementApplication/PostDoctoralManagementApplication-ejb/src/main/setup/glassfish-resources.xml}}) most be loaded into the glassfish server by going to the Resources (under common task) - Add Resources (button) and follow prompts.
\textbf{Note:} The @user and @password of the JDBC resources needs to match up with the ones used in the MySQL server. Also note that glassfish server runs using the java runtime environment and if there multiple version of the environments on a machine it might not run on your desired one. So caution must be taking when running the glassfish server for deployment.\\
\subsection{MySQL 5.6}
MySQL 5.6 acts as the persistance layer implementation of the system.
\subsubsection{Installation}
All the required installation knowledge can be found on this link \textbf{\seqsplit{http://dev.mysql.com/doc/refman/5.6/en/installing.html}}
\subsubsection{Configuration}
A @user needs to be setup for which all the required databases will be accessed under. This user information needs to be propagated to the glassfish resources in both the configured glassfish server and the projects glassfish-resourcs.xml file.
\\
\\
To setup the tables required by the application the initilizing script (found under ~/Source\_Code/Database/InitialisePostDoc\_DB) needs to be in the configured MySQL. This action can be accomplished through running the script on MySQL workbench (which is avaible on both operating systems and should be seen in the availible via the installation).
\subsection{Netbeans}
Netbeans is the primary IDE in which the application was developed under.
\subsubsection{Installation}
All required information for installation can be found here \textbf{\seqsplit{https://netbeans.org/community/releases/80/install.html}}.
\subsubsection{Configuration}
The IDE must have the configured glassfish server registered under the services in its services group. All required information for this task can be found here \textbf{\seqsplit{http://docs.oracle.com/cd/E19798-01/821-1770/gioew/index.html}}
\\
The IDE must have Java 8 sdk as the default Java platform.\\
\\
At this point all the required software should have been installed ann properly configured. Now we shall focus on the actual deployment of the document.
\section{Deployment}
Glassfish server provideds the deployment methods which we are going to use to achieve both production and development deployment. The two forms of deployment could be achieved by the following steps:
\subsection{Production Deployment}
\subsubsection{Autodeployment}
Glassfish provides a number of ways to deploy an application. Both are provided under the assumption the your glassfish server runs using the java 8 runtime environment. The autodeploy method requires that the .war file be situated in the autodeploy (\textbf{\seqsplit{~/glassfish-4.0/glassfish/domains/domain1/autodeploy/}}) folder.\\
\\
Hence, PostDoctoralManagementApplication-web-0.0.war (found under \textbf{\seqsplit{~/PostDoctoralManagementApplication/PostDoctoralManagementApplication-ear/target/PostDoctoralManagementApplication-ear-0.0/}}) should be copied to the autodeploy folder. The application should be ran when the server starts.
\subsubsection{On demand deployment}
The user is required to be logged on to the running glassfish server admin page (at the web address localhost:4848). From there the on demand application deployment should be on the home page of the adminstration.
\subsection{Developer Deployment}
A developer could deploy the app using Netbeans "run web application" option (which will start the glassfish server and deploy on it).\\
\\
All these steps will provide the needed actions to successfully deploy the application. The variants between windows and linux lay within the installions from which point both processes are similar.
\end{document}
